% ===================================================================
%  4. ANALISI DEI REQUISITI E TRACCIABILITÀ
% ===================================================================
\newpage
\section{Analisi dei Requisiti}

\subsection{Introduzione e Classificazione}
In questa sezione vengono elencati, classificati e descritti formalmente tutti i requisiti individuati durante la fase di analisi del dominio, dallo studio del capitolato d'appalto e dalle decisioni architetturali interne. I requisiti rappresentano i vincoli operativi e strutturali che il sistema \textbf{SmartOrder} deve soddisfare.

I requisiti sono identificati da un codice univoco secondo la seguente struttura:
\begin{center}
    \textbf{[TIPO]-[PRIORITÀ]\_[ID]}
\end{center}

\textbf{TIPO di requisito:}
\begin{itemize}
    \item \textbf{RF}: Requisito Funzionale (Cosa fa il sistema).
    \item \textbf{RQA}: Requisito Qualitativo / Non Funzionale (Come si comporta il sistema).
    \item \textbf{RP}: Requisito Prestazionale (Limiti e tempi di risposta).
    \item \textbf{RV}: Requisito di Vincolo (Scelte tecnologiche e architetturali obbligate).
\end{itemize}

\textbf{PRIORITÀ del requisito:}
\begin{itemize}
    \item \textbf{OB (Obbligatorio)}: Requisito fondamentale per il funzionamento base.
    \item \textbf{DE (Desiderabile)}: Requisito a valore aggiunto, utile ma non strettamente vitale.
    \item \textbf{OP (Opzionale)}: Possibile sviluppo futuro, non previsto in questa fase.
\end{itemize}

% -------------------------------------------------------------------
\newpage
\subsection{Requisiti Funzionali}

\subsubsection{Requisiti Funzionali obbligatori}
\renewcommand{\arraystretch}{1.15}
\begin{longtable}{|>{\raggedright}p{0.18\textwidth}|>{\raggedright\arraybackslash}p{0.5\textwidth}|>{\raggedright\arraybackslash}p{0.2\textwidth}|}

    \hline
    \rowcolor[gray]{0.9}
    \textbf{Codice} & \textbf{Descrizione} & \textbf{Fonti} \\
    \hline
    \endfirsthead
    
    \hline
    \rowcolor[gray]{0.9}
    \textbf{Codice} & \textbf{Descrizione} & \textbf{Fonti}  \\
    \hline
    \endhead

    \hline
    RF-OB\_01 \hypertarget{rf-ob_1} & Il sistema deve permettere all'Utente di inserire il proprio indirizzo email per fare il login. & \hyperref[uc_01.1]{UC\_01.1} \\
    \hline
    RF-OB\_02 \hypertarget{rf-ob_2} & L'Utente deve poter inserire la propria email nel sistema. & \hyperref[uc_01.1]{UC\_01.1} \\
    \hline
    RF-OB\_03 \hypertarget{rf-ob_3} & Il sistema deve permettere all'Utente di inserire la password in modo oscurato (nascosta per sicurezza). & \hyperref[uc_01.2]{UC\_01.2} \\
    \hline
    RF-OB\_04 \hypertarget{rf-ob_4} & L'Utente deve poter inserire la propria password a sistema. & \hyperref[uc_01.2]{UC\_01.2} \\
    \hline
    RF-OB\_05 \hypertarget{rf-ob_5} & Il sistema deve fornire un pulsante o comando per confermare l'accesso. & \hyperref[uc_01]{UC\_01} \\
    \hline
    RF-OB\_06 \hypertarget{rf-ob_6} & L'Utente deve poter inviare le credenziali inserite per la validazione. & \hyperref[uc_01]{UC\_01} \\
    \hline
    RF-OB\_07 \hypertarget{rf-ob_7} & Il sistema deve mostrare un errore se l'email inserita non ha un formato valido. & \hyperref[uc_01]{UC\_01} \\
    \hline
    RF-OB\_08 \hypertarget{rf-ob_8} & Il sistema deve mostrare un errore se l'email inserita non è presente nel database locale. & \hyperref[uc_01]{UC\_01} \\
    \hline
    RF-OB\_09 \hypertarget{rf-ob_9} & Il sistema deve mostrare un errore se la combinazione di email e password è sbagliata. & \hyperref[uc_01]{UC\_01} \\
    \hline
    RF-OB\_10 \hypertarget{rf-ob_10} & Dopo un login valido, il sistema deve creare la sessione e far accedere l'utente alla propria area (Cliente o Operatore). & \hyperref[uc_01]{UC\_01} \\
    \hline
    RF-OB\_11 \hypertarget{rf-ob_11} & Il sistema deve fornire un pulsante per il logout sicuro. & \hyperref[uc_03]{UC\_03} \\
    \hline
    RF-OB\_12 \hypertarget{rf-ob_12} & L'Utente deve poter chiudere la propria sessione sul server. & \hyperref[uc_03]{UC\_03} \\
    \hline
    RF-OB\_13 \hypertarget{rf-ob_13} & Il sistema deve mostrare i dati aziendali del Cliente (ID e Ragione Sociale) in sola lettura. & \hyperref[uc_00.2]{UC\_00.2} \\
    \hline
    RF-OB\_14 \hypertarget{rf-ob_14} & Il sistema deve mostrare all'Utente il proprio indirizzo email in sola lettura. & \hyperref[uc_04.1]{UC\_04.1} \\
    \hline
    RF-OB\_15 \hypertarget{rf-ob_15} & Il sistema deve mostrare in tempo reale tutti i messaggi inviati dall'Utente nella chat corrente. & \hyperref[uc_07]{UC\_07} \\
    \hline
    RF-OB\_16 \hypertarget{rf-ob_16} & Il sistema deve mostrare in tempo reale le risposte testuali scritte dall'AI o dall'Operatore di supporto. & \hyperref[uc_07]{UC\_07} \\
    \hline
    RF-OB\_17 \hypertarget{rf-ob_17} & Il sistema deve avere un campo di testo per l'inserimento dei messaggi. & \hyperref[uc_08.1]{UC\_08.1} \\
    \hline
    RF-OB\_18 \hypertarget{rf-ob_18} & L'Utente deve poter scrivere e inviare il suo ordine tramite messaggio di testo. & \hyperref[uc_08.1]{UC\_08.1} \\
    \hline
    RF-OB\_19 \hypertarget{rf-ob_19} & Il sistema deve bloccare l'invio e mostrare un errore se il messaggio di testo supera il limite di caratteri. & \hyperref[uc_08.1]{UC\_08.1} \\
    \hline
    RF-OB\_20 \hypertarget{rf-ob_20} & Il sistema deve permettere la registrazione e l'invio di messaggi vocali. & \hyperref[uc_08.2]{UC\_08.2} \\
    \hline
    RF-OB\_21 \hypertarget{rf-ob_21} & L'Utente deve poter usare il microfono per registrare e inviare l'ordine a voce. & \hyperref[uc_08.2]{UC\_08.2} \\
    \hline
    RF-OB\_22 \hypertarget{rf-ob_22} & Il sistema deve convertire l'audio in testo usando un modulo Speech-to-Text. & \hyperref[uc_08.2]{UC\_08.2} \\
    \hline
    RF-OB\_23 \hypertarget{rf-ob_23} & Il sistema deve bloccare l'invio e mostrare un errore se il file audio pesa più di 10MB. & \hyperref[uc_08.2]{UC\_08.2}, Capitolato \\
    \hline
    RF-OB\_24 \hypertarget{rf-ob_24} & Il Cliente deve poter resettare la chat con l'AI e svuotare il carrello attuale. & \hyperref[uc_10]{UC\_10} \\
    \hline
    RF-OB\_25 \hypertarget{rf-ob_25} & L'AI (LLM) deve capire quali prodotti vuole l'utente e in che quantità leggendo/ascoltando il messaggio. & Capitolato \\
    \hline
    RF-OB\_26 \hypertarget{rf-ob_26} & Il sistema deve cercare i prodotti nel database usando la ricerca vettoriale (Vector Search). & Interno \\
    \hline
    RF-OB\_27 \hypertarget{rf-ob_27} & Il sistema deve aggiungere in automatico i prodotti trovati al carrello se l'AI è sicura della sua scelta per oltre il 70\%. & Interno \\
    \hline
    RF-OB\_28 \hypertarget{rf-ob_28} & Vincolo Assortimento: Il sistema deve ignorare i prodotti che non fanno parte del catalogo assegnato a quel Cliente (tabella \texttt{asscli}). & Interno \\
    \hline
    RF-OB\_29 \hypertarget{rf-ob_29} & Vincolo Stato: Il sistema deve impedire l'aggiunta di prodotti segnati come "non disponibili" o "sospesi" nel database (\texttt{anaart}). & Interno \\
    \hline
    RF-OB\_30 \hypertarget{rf-ob_30} & Deduzione Formati: L'AI deve convertire automaticamente le quantità richieste nelle giuste unità di vendita (es. da bottiglie a casse). & Interno \\
    \hline
    RF-OB\_31 \hypertarget{rf-ob_31} & Il sistema deve fornire una barra di ricerca per cercare manualmente i prodotti nel catalogo. & \hyperref[uc_11]{UC\_11} \\
    \hline
    RF-OB\_32 \hypertarget{rf-ob_32} & L'Utente deve poter inserire del testo per cercare un prodotto. & \hyperref[uc_11]{UC\_11} \\
    \hline
    RF-OB\_33 \hypertarget{rf-ob_33} & Il sistema deve mostrare i risultati della ricerca in tempo reale (Live Search). & \hyperref[uc_11]{UC\_11} \\
    \hline
    RF-OB\_34 \hypertarget{rf-ob_34} & Il sistema deve mostrare per ogni prodotto trovato: Nome, ID, Linea, Famiglia e Modalità di Vendita. & \hyperref[uc_00.1]{UC\_00.1} \\
    \hline
    RF-OB\_35 \hypertarget{rf-ob_35} & L'Utente deve poter cliccare su un risultato per decidere quanti pezzi acquistarne. & \hyperref[uc_12]{UC\_12} \\
    \hline
    RF-OB\_36 \hypertarget{rf-ob_36} & L'Utente deve poter modificare la quantità usando i tasti +/- o scrivendo direttamente il numero. & \hyperref[uc_12]{UC\_12} \\
    \hline
    RF-OB\_37 \hypertarget{rf-ob_37} & Il sistema deve correggere in automatico la quantità se l'utente inserisce numeri negativi, zero o troppo grandi. & \hyperref[uc_12]{UC\_12} \\
    \hline
    RF-OB\_38 \hypertarget{rf-ob_38} & L'Utente deve poter confermare l'aggiunta del prodotto al carrello della sessione corrente. & \hyperref[uc_12]{UC\_12} \\
    \hline
    RF-OB\_39 \hypertarget{rf-ob_39} & Il sistema deve mostrare un avviso di "Carrello Vuoto" se non ci sono prodotti inseriti. & \hyperref[uc_13]{UC\_13} \\
    \hline
    RF-OB\_40 \hypertarget{rf-ob_40} & Il sistema deve mostrare i dettagli logistici e le quantità di tutti i prodotti presenti nel carrello. & \hyperref[uc_13]{UC\_13} \\
    \hline
    RF-OB\_41 \hypertarget{rf-ob_41} & L'Utente deve poter sbloccare la riga di un prodotto nel carrello per poterne modificare la quantità. & \hyperref[uc_14.2]{UC\_14.2} \\
    \hline
    RF-OB\_42 \hypertarget{rf-ob_42} & L'Utente deve poter salvare la nuova quantità modificata per un prodotto nel carrello. & \hyperref[uc_14]{UC\_14} \\
    \hline
    RF-OB\_43 \hypertarget{rf-ob_43} & L'Utente deve poter rimuovere definitivamente un articolo dal carrello. & \hyperref[uc_14.1]{UC\_14.1} \\
    \hline
    RF-OB\_44 \hypertarget{rf-ob_44} & Il sistema deve disabilitare il pulsante per inviare l'ordine se il carrello è vuoto. & \hyperref[uc_15]{UC\_15} \\
    \hline
    RF-OB\_45 \hypertarget{rf-ob_45} & Il Cliente deve poter chiudere il carrello e inviare l'ordine definitivo. & \hyperref[uc_15]{UC\_15} \\
    \hline
    RF-OB\_46 \hypertarget{rf-ob_46} & Il sistema deve bloccare l'invio dell'ordine e segnalare gli errori se ci sono problemi con le regole del gestionale (es. prodotto fuori catalogo). & \hyperref[uc_15]{UC\_15} \\
    \hline
    RF-OB\_47 \hypertarget{rf-ob_47} & Dopo l'invio dell'ordine, il sistema deve svuotare la sessione e reindirizzare il Cliente alla pagina dello Storico Ordini. & \hyperref[uc_15]{UC\_15}, Interno \\
    \hline
    RF-OB\_48 \hypertarget{rf-ob_48} & Il sistema deve creare un file JSON formattato contenente tutti i dati dell'ordine appena chiuso. & \hyperref[uc_15]{UC\_15}, Capitolato \\
    \hline
    RF-OB\_49 \hypertarget{rf-ob_49} & Il Cliente deve poter visualizzare la lista dei propri ordini passati, ordinati per data. & \hyperref[uc_16]{UC\_16} \\
    \hline
    RF-OB\_50 \hypertarget{rf-ob_50} & Il sistema deve mostrare per ogni ordine storico: ID, Data, totale degli articoli e una piccola anteprima dei prodotti. & \hyperref[uc_00.3]{UC\_00.3}, \hyperref[uc_16]{UC\_16} \\
    \hline
    RF-OB\_51 \hypertarget{rf-ob_51} & Il Cliente deve poter cliccare su un ordine passato per vederne tutti i dettagli logistici. & \hyperref[uc_18]{UC\_18} \\
    \hline
    RF-OB\_52 \hypertarget{rf-ob_52} & Il sistema deve mostrare in sola lettura la chat originale tra utente e AI che ha generato quell'ordine specifico. & \hyperref[uc_00.4]{UC\_00.4} \\
    \hline
    RF-OB\_53 \hypertarget{rf-ob_53} & Il Cliente deve poter richiedere l'aiuto di un operatore umano, mettendo in pausa l'AI. & \hyperref[uc_19]{UC\_19} \\
    \hline
    RF-OB\_54 \hypertarget{rf-ob_54} & L'Operatore deve poter accedere alla dashboard con la lista dei Ticket Aperti (richieste d'aiuto). & \hyperref[uc_21]{UC\_21} \\
    \hline
    RF-OB\_55 \hypertarget{rf-ob_55} & Nella lista dei ticket, il sistema deve mostrare da quanto tempo il Cliente sta aspettando una risposta. & \hyperref[uc_21.5]{UC\_21.5} \\
    \hline
    RF-OB\_56 \hypertarget{rf-ob_56} & Il sistema deve aggiornare in tempo reale (Live Sync) la lista dei ticket dell'Operatore all'arrivo di nuove richieste, senza ricaricare la pagina. & \hyperref[uc_21]{UC\_21} \\
    \hline
    RF-OB\_57 \hypertarget{rf-ob_57} & L'Operatore deve poter prendere in carico un ticket bloccandolo per gli altri colleghi (Lock) per aprire la console di assistenza. & \hyperref[uc_22]{UC\_22} \\
    \hline
    RF-OB\_58 \hypertarget{rf-ob_58} & L'Operatore deve poter scrivere e inviare messaggi in chat direttamente al Cliente per fornirgli assistenza. & \hyperref[uc_08.1]{UC\_08.1} \\
    \hline
    RF-OB\_59 \hypertarget{rf-ob_59} & L'Operatore deve poter usare la ricerca per aggiungere forzatamente prodotti al carrello del Cliente che sta aiutando. & \hyperref[uc_11]{UC\_11}, \hyperref[uc_12]{UC\_12} \\
    \hline
    RF-OB\_60 \hypertarget{rf-ob_60} & L'Operatore deve poter modificare le quantità o rimuovere prodotti dal carrello del Cliente. & \hyperref[uc_13]{UC\_13}, \hyperref[uc_14]{UC\_14} \\
    \hline
    RF-OB\_61 \hypertarget{rf-ob_61} & L'Operatore deve poter chiudere il ticket per segnare il problema come risolto e sbloccare la sessione (rimuovere il Lock). & \hyperref[uc_22.2]{UC\_22.2} \\
    \hline
    RF-OB\_62 \hypertarget{rf-ob_62} & Il Cliente deve poter chiudere il ticket da solo se risolve il problema o non ha più bisogno di aiuto. & \hyperref[uc_20]{UC\_20} \\
    \hline
    RF-OB\_63 \hypertarget{rf-ob_63} & L'Operatore deve poter accedere alla dashboard che mostra tutti gli ordini inviati da tutti i Clienti. & \hyperref[uc_23]{UC\_23} \\
    \hline
    RF-OB\_64 \hypertarget{rf-ob_64} & L'Operatore deve poter ordinare le tabelle in modo crescente o decrescente cliccando sui nomi delle colonne. & \hyperref[uc_21]{UC\_21}, \hyperref[uc_23]{UC\_23}, \hyperref[uc_25]{UC\_25} \\
    \hline
    RF-OB\_65 \hypertarget{rf-ob_65} & L'Operatore deve poter cliccare su un ordine di qualsiasi cliente per aprirlo e vederne i dettagli. & \hyperref[uc_24]{UC\_24} \\
    \hline
    RF-OB\_66 \hypertarget{rf-ob_66} & L'Operatore deve poter filtrare i dati delle tabelle scegliendo una data di inizio e di fine. & \hyperref[uc_21]{UC\_21}, \hyperref[uc_23]{UC\_23}, \hyperref[uc_25]{UC\_25} \\
    \hline
    RF-OB\_67 \hypertarget{rf-ob_67} & L'Operatore deve poter cercare del testo solo all'interno di una colonna specifica (es. cercare solo l'ID Ordine). & \hyperref[uc_21]{UC\_21}, \hyperref[uc_23]{UC\_23} \\
    \hline
    RF-OB\_68 \hypertarget{rf-ob_68} & L'Operatore deve poter usare una barra di ricerca per cercare testo liberamente in tutte le colonne della tabella (Ricerca Full-Text). & \hyperref[uc_25]{UC\_25} \\
    \hline
    RF-OB\_69 \hypertarget{rf-ob_69} & Il sistema deve salvare nel database tutto lo storico delle chat e delle decisioni prese dall'AI. & Interno, Capitolato \\
    \hline
    RF-OB\_70 \hypertarget{rf-ob_70} & Il sistema deve mostrare la percentuale di sicurezza dell'AI (Confidence Score) solo nelle schermate dedicate all'Operatore. & \hyperref[uc_22.1]{UC\_22.1}, \hyperref[uc_24.1]{UC\_24.1} \\
    \hline
    RF-OB\_71 \hypertarget{rf-ob_71} & Il sistema deve mostrare schermate di caricamento (spinner o barre di avanzamento) durante le operazioni che richiedono elaborazioni di lunga durata. & Interno \\
    \hline
    RF-OB\_72 \hypertarget{rf-ob_72} & Le API del server (Backend) devono rifiutare qualsiasi richiesta priva di un token di accesso valido per la sessione corrente. & Interno \\
    \hline
    RF-OB\_73 \hypertarget{rf-ob_73} & Il sistema deve caricare i dati delle tabelle in modo incrementale (es. Paginazione o Lazy Loading) per evitare il blocco dell'interfaccia in presenza di grandi volumi di dati. & Interno \\
    \hline
    RF-OB\_74 \hypertarget{rf-ob_74} & Gli utenti non possono registrarsi autonomamente sull'applicazione: gli account vengono creati e gestiti esclusivamente all'interno del gestionale ERP aziendale. & Interno (Logica B2B) \\
    \hline
    RF-OB\_75 \hypertarget{rf-ob_75} & L'applicazione web non deve eseguire calcoli di sconti, prezzi o totali in euro: tutta la logica di fatturazione è delegata esclusivamente al gestionale ERP. & Interno (Logica ERP) \\
    \hline
    \multicolumn{3}{c}{} \\
    \caption{Requisiti Funzionali Obbligatori}\hypertarget{tab:funzionali_obbligatori}\\
\end{longtable}

    
\subsubsection{Requisiti Funzionali desiderabili}
\renewcommand{\arraystretch}{1.15}
\begin{longtable}{|>{\raggedright}p{0.18\textwidth}|>{\raggedright\arraybackslash}p{0.5\textwidth}|>{\raggedright\arraybackslash}p{0.2\textwidth}|}
    \hline
    \rowcolor[gray]{0.9}
    \textbf{Codice} & \textbf{Descrizione} & \textbf{Fonti} \\
    \hline
    \endfirsthead
    
    \hline
    \rowcolor[gray]{0.9}
    \textbf{Codice} & \textbf{Descrizione} & \textbf{Fonti}  \\
    \hline
    \endhead
    
    \hline
    RF-DE\_01 \hypertarget{rf-de_1} & Il sistema deve fornire un link per accedere alla pagina di recupero della password dimenticata. & \hyperref[uc_02]{UC\_02} \\
    \hline
    RF-DE\_02 \hypertarget{rf-de_2} & L'Utente deve poter richiedere il reset della password inserendo la propria email. & \hyperref[uc_02]{UC\_02}, \hyperref[uc_02.1]{UC\_02.1} \\
    \hline
    RF-DE\_03 \hypertarget{rf-de_3} & Il sistema deve permettere all'Utente loggato di avviare la procedura per cambiare la propria password. & \hyperref[uc_04.2]{UC\_04.2} \\
    \hline
    RF-DE\_04 \hypertarget{rf-de_4} & L'Utente deve poter inviare la vecchia password e la nuova password scelta per salvarla nel sistema. & \hyperref[uc_05.1]{UC\_05.1}, \hyperref[uc_05.2]{UC\_05.2}, \hyperref[uc_05.3]{UC\_05.3} \\
    \hline
    RF-DE\_05 \hypertarget{rf-de_5} & Il sistema deve bloccare il cambio password se la vecchia password è errata o se le due nuove password inserite non coincidono. & \hyperref[uc_05.1]{UC\_05.1}, \hyperref[uc_05.3]{UC\_05.3} \\
    \hline
    RF-DE\_06 \hypertarget{rf-de_6} & Il sistema deve permettere l'invio di file immagine all'interno della chat. & \hyperref[uc_08.3]{UC\_08.3} \\
    \hline
    RF-DE\_07 \hypertarget{rf-de_7} & L'Utente deve poter inviare un'immagine scattando una foto o scegliendola dalla galleria del proprio dispositivo. & \hyperref[uc_08.3]{UC\_08.3} \\
    \hline
    RF-DE\_08 \hypertarget{rf-de_8} & Il sistema deve bloccare l'invio di immagini e mostrare un errore se il file pesa più di 5MB. & \hyperref[uc_08.3]{UC\_08.3}, Interno \\
    \hline
    RF-DE\_09 \hypertarget{rf-de_9} & Il Cliente deve poter valutare la qualità delle risposte dell'AI lasciando un feedback (es. pollice su / pollice giù). & \hyperref[uc_09.1]{UC\_09.1}, \hyperref[uc_09.2]{UC\_09.2} \\
    \hline
    RF-DE\_10 \hypertarget{rf-de_10} & Se il Cliente dà un voto negativo all'AI, il sistema deve richiedere che indichi il tipo di errore (tramite tag predisposti). & \hyperref[uc_09.2]{UC\_09.2} \\
    \hline
    RF-DE\_11 \hypertarget{rf-de_11} & Il Cliente deve poter aggiungere un testo libero per spiegare meglio il motivo dell'errore segnalato. & \hyperref[uc_09.2]{UC\_09.2} \\
    \hline
    RF-DE\_12 \hypertarget{rf-de_12} & Il Cliente deve poter scaricare l'elenco dei propri ordini storici in un file formato CSV. & \hyperref[uc_17]{UC\_17} \\
    \hline
    RF-DE\_13 \hypertarget{rf-de_13} & L'Operatore deve poter accedere alla dashboard per il controllo degli errori e l'addestramento dell'AI (Ottimizzazione AI). & \hyperref[uc_25]{UC\_25} \\
    \hline
    RF-DE\_14 \hypertarget{rf-de_14} & Il sistema deve mostrare all'Operatore la lista degli errori segnalati dai clienti e le correzioni manuali fatte ai carrelli. & \hyperref[uc_25]{UC\_25} \\
    \hline
    RF-DE\_15 \hypertarget{rf-de_15} & L'Operatore deve poter filtrare le segnalazioni di errore in base al loro stato (es. "Da gestire", "Accettate"). & \hyperref[uc_25]{UC\_25} \\
    \hline
    RF-DE\_16 \hypertarget{rf-de_16} & L'Operatore deve poter aprire un singolo errore AI per rileggere tutta la chat originale tra il Cliente e il bot. & \hyperref[uc_26]{UC\_26} \\
    \hline
    RF-DE\_17 \hypertarget{rf-de_17} & L'Operatore deve poter approvare una segnalazione di errore, confermando che i dati verranno usati per riaddestrare l'AI. & \hyperref[uc_27.1]{UC\_27.1} \\
    \hline
    RF-DE\_18 \hypertarget{rf-de_18} & L'Operatore deve poter scartare una segnalazione se non si tratta di un vero errore dell'AI (es. l'utente ha scritto male). & \hyperref[uc_27.2]{UC\_27.2} \\
    \hline
    RF-DE\_19 \hypertarget{rf-de_19} & L'Operatore deve poter annullare l'approvazione di una segnalazione se l'aveva confermata per sbaglio (Rollback). & \hyperref[uc_28]{UC\_28} \\
    \hline
    \multicolumn{3}{c}{} \\
    \caption{Requisiti Funzionali Desiderabili}\hypertarget{tab:funzionali_desiderabili}\\
\end{longtable}

    
\subsubsection{Requisiti Funzionali opzionali}
\renewcommand{\arraystretch}{1.15}
\begin{longtable}{|>{\raggedright}p{0.18\textwidth}|>{\raggedright\arraybackslash}p{0.5\textwidth}|>{\raggedright\arraybackslash}p{0.2\textwidth}|}
    \hline
    \rowcolor[gray]{0.9}
    \textbf{Codice} & \textbf{Descrizione} & \textbf{Fonti} \\
    \hline
    \endfirsthead
    
    \hline
    \rowcolor[gray]{0.9}
    \textbf{Codice} & \textbf{Descrizione} & \textbf{Fonti}  \\
    \hline
    \endhead
    
    \hline
    RF-OP\_01 \hypertarget{rf-op_1} & Il sistema deve permettere l'invio di messaggi d'ordine e la loro elaborazione tramite l'app di WhatsApp. & Interno \\
    \hline
    RF-OP\_02 \hypertarget{rf-op_2} & Il sistema deve permettere l'invio di messaggi d'ordine e la loro elaborazione tramite l'app di Telegram. & Interno \\
    \hline
    RF-OP\_03 \hypertarget{rf-op_3} & Il sistema deve generare grafici e statistiche generali sull'uso della piattaforma, mostrandoli all'Operatore. & Interno \\
    \hline
    \multicolumn{3}{c}{} \\
    \caption{Requisiti Funzionali Opzionali}\hypertarget{tab:funzionali_opzionali}\\

\end{longtable}


% -------------------------------------------------------------------
\newpage
\subsection{Requisiti Qualitativi (Non Funzionali)}
Questa sezione traccia i requisiti legati alle qualità del software, quali usabilità, sicurezza e manutenzione (identificati dal prefisso \textbf{RQA}).

\subsubsection{Requisiti Qualitativi obbligatori}
\renewcommand{\arraystretch}{1.15}
\begin{longtable}{|>{\raggedright}p{0.18\textwidth}|>{\raggedright\arraybackslash}p{0.5\textwidth}|>{\raggedright\arraybackslash}p{0.2\textwidth}|}

    \hline
    \rowcolor[gray]{0.9}
    \textbf{Codice} & \textbf{Descrizione} & \textbf{Fonti} \\
    \hline
    \endfirsthead
    
    \hline
    \rowcolor[gray]{0.9}
    \textbf{Codice} & \textbf{Descrizione} & \textbf{Fonti}  \\
    \hline
    \endhead

    \hline
    RQA-OB\_01 \hypertarget{rqa-ob_1} & L'interfaccia del Cliente deve essere progettata con approccio \textit{Mobile-first}: tutti i controlli interattivi devono avere un'area toccabile di almeno 44$\times$44 pixel e l'interfaccia non deve presentare scroll orizzontale su alcun viewport di larghezza supportata. & Interno \\
    \hline
    RQA-OB\_02 \hypertarget{rqa-ob_2} & L'interfaccia dell'Operatore (Backoffice) deve essere progettata con approccio \textit{Desktop-first}: tutti i pannelli e i controlli devono essere completamente fruibili su viewport di larghezza non inferiore a 1024 pixel, senza degradazioni di layout. & Interno \\
    \hline
    RQA-OB\_03 \hypertarget{rqa-ob_3} & Il codice del sistema deve essere salvato su Git e documentato seguendo le regole scritte nel documento delle \textit{Norme di Progetto}. & Norme di Progetto \\
    \hline
    \multicolumn{3}{c}{} \\
    \caption{Requisiti Qualitativi obbligatori}\hypertarget{tab:qualitativi_obbligatori}\\
\end{longtable}


% -------------------------------------------------------------------
\newpage
\subsection{Requisiti Prestazionali}
Questa sezione raccoglie i limiti tecnici e i tempi di risposta che il sistema deve rispettare (identificati dal prefisso \textbf{RP}).

\subsubsection{Requisiti Prestazionali obbligatori}
\renewcommand{\arraystretch}{1.15}
\begin{longtable}{|>{\raggedright}p{0.18\textwidth}|>{\raggedright\arraybackslash}p{0.5\textwidth}|>{\raggedright\arraybackslash}p{0.2\textwidth}|}

    \hline
    \rowcolor[gray]{0.9}
    \textbf{Codice} & \textbf{Descrizione} & \textbf{Fonti} \\
    \hline
    \endfirsthead
    
    \hline
    \rowcolor[gray]{0.9}
    \textbf{Codice} & \textbf{Descrizione} & \textbf{Fonti}  \\
    \hline
    \endhead

    \hline
    RP-OB\_01 \hypertarget{rp-ob_1} & Il modello AI (LLM) deve completare l'elaborazione del messaggio e rispondere entro massimo 20 secondi, misurati come valore mediano su un campione di almeno 50 richieste consecutive in assenza di traffico concorrente. & Interno, Capitolato \\
    \hline
    RP-OB\_02 \hypertarget{rp-ob_2} & Il sistema di trascrizione vocale (Speech-to-Text) deve riconoscere le parole con un Word Error Rate (WER) inferiore al 10\%, misurato su un campione di almeno 100 registrazioni audio in italiano standard. & Interno \\
    \hline
    RP-OB\_03 \hypertarget{rp-ob_3} & La percentuale di sicurezza (Confidence Score) calcolata dall'AI deve superare il 70\% per poter aggiungere prodotti in automatico, misurata su un campione di almeno 200 richieste di test con prodotti presenti nel catalogo corrente. & Interno \\
    \hline
    RP-OB\_04 \hypertarget{rp-ob_4} & L'applicazione web deve caricare gli elementi principali a schermo in meno di 3 secondi, misurati al 90° percentile su 10 caricamenti consecutivi da una connessione con banda minima di 20\,Mbit/s in download. & Interno \\
    \hline
    RP-OB\_05 \hypertarget{rp-ob_5} & Quando l'utente usa la barra di ricerca prodotti, i risultati devono comparire in meno di 500 millisecondi, misurati al 95° percentile su un database contenente almeno 2000 prodotti e con un massimo di 10 utenti concorrenti. & Interno \\
    \hline
    \multicolumn{3}{c}{} \\
    \caption{Requisiti Prestazionali obbligatori}\hypertarget{tab:prestazionali_obbligatori}\\
\end{longtable}

\subsubsection{Requisiti Prestazionali desiderabili}
\renewcommand{\arraystretch}{1.15}
\begin{longtable}{|>{\raggedright}p{0.18\textwidth}|>{\raggedright\arraybackslash}p{0.5\textwidth}|>{\raggedright\arraybackslash}p{0.2\textwidth}|}

    \hline
    \rowcolor[gray]{0.9}
    \textbf{Codice} & \textbf{Descrizione} & \textbf{Fonti} \\
    \hline
    \endfirsthead
    
    \hline
    \rowcolor[gray]{0.9}
    \textbf{Codice} & \textbf{Descrizione} & \textbf{Fonti}  \\
    \hline
    \endhead

    \hline
    RP-DE\_01 \hypertarget{rp-de_1} & L'estrazione e il download in CSV dello storico ordini deve concludersi in massimo 5 secondi per liste con meno di 1000 ordini, misurato al 90° percentile su 10 esecuzioni consecutive in condizioni di rete locale (LAN). & Interno \\
    \hline
    \multicolumn{3}{c}{} \\
    \caption{Requisiti Prestazionali desiderabili}\hypertarget{tab:prestazionali_desiderabili}\\
\end{longtable}


% -------------------------------------------------------------------
\newpage
\subsection{Requisiti di Vincolo}
Questa sezione elenca gli obblighi e le scelte tecnologiche imposte dall'esterno a cui il sistema deve conformarsi (identificati dal prefisso \textbf{RV}).

\subsubsection{Requisiti di Vincolo obbligatori}
\renewcommand{\arraystretch}{1.15}
\begin{longtable}{|>{\raggedright}p{0.18\textwidth}|>{\raggedright\arraybackslash}p{0.5\textwidth}|>{\raggedright\arraybackslash}p{0.2\textwidth}|}

    \hline
    \rowcolor[gray]{0.9}
    \textbf{Codice} & \textbf{Descrizione} & \textbf{Fonti} \\
    \hline
    \endfirsthead
    
    \hline
    \rowcolor[gray]{0.9}
    \textbf{Codice} & \textbf{Descrizione} & \textbf{Fonti}  \\
    \hline
    \endhead

    \hline
    RV-OB\_01 \hypertarget{rv-ob_1} & Il Frontend (la parte grafica usata dall'utente) e il Backend (il server) devono essere programmati come due sistemi separati che comunicano tra loro. & Interno, Norme di Progetto \\
    \hline
    RV-OB\_02 \hypertarget{rv-ob_2} & Il database della piattaforma web deve essere di tipo relazionale, nello specifico PostgreSQL. & Interno \\
    \hline
    RV-OB\_03 \hypertarget{rv-ob_3} & La comprensione del testo scritto dagli utenti deve essere gestita tramite tecnologie AI, nello specifico gli Embedding e la Vector Search. & Capitolato \\
    \hline
    RV-OB\_04 \hypertarget{rv-ob_4} & Il file JSON dell'ordine deve essere posizionato in una cartella o sistema condiviso in modo che il gestionale aziendale (ERP) possa leggerlo. & Interno (Logica ERP) \\
    \hline
    RV-OB\_05 \hypertarget{rv-ob_5} & Tutte le comunicazioni tra la webapp e il server devono usare esclusivamente connessioni sicure crittografate (HTTPS). & Interno \\
    \hline
    RV-OB\_06 \hypertarget{rv-ob_6} & Le password non devono mai essere salvate in chiaro, ma protette nel database usando l'algoritmo di hashing scrypt e salt. & Interno \\
    \hline
    RV-OB\_07 \hypertarget{rv-ob_7} & Il salvataggio nel database e l'uso delle email degli utenti deve rispettare le regole europee sulla privacy (GDPR). & Interno \\
    \hline
    RV-OB\_08 \hypertarget{rv-ob_8} & L'applicazione web deve funzionare correttamente sulle ultime due versioni stabili dei principali browser desktop e mobile: Google Chrome, Mozilla Firefox, Apple Safari e Microsoft Edge. & Interno \\
    \hline
    RV-OB\_09 \hypertarget{rv-ob_9} & L'interfaccia del Cliente deve adattarsi graficamente a schermi stretti, con una larghezza minima di 375 pixel. & Interno \\
    \hline
    RV-OB\_10 \hypertarget{rv-ob_10} & L'interfaccia dell'Operatore è progettata per essere usata su monitor da computer, richiedendo uno schermo largo almeno 1024 pixel. & Interno \\
    \hline
    \multicolumn{3}{c}{} \\
    \caption{Requisiti di Vincolo obbligatori}\hypertarget{tab:vincolo_obbligatori}\\
\end{longtable}